% 原书第 10.4 和 10.5 节。
\section{新的前沿}
\label{sec:ch10-new-frontiers}

\subsection{大语言模型智能体}
\label{sec:ch10-agents}
大语言模型智能体是一种利用大语言模型完成复杂任务的框架。此类智能体是模块化程序，可以读取用户请求，推理完成请求所需的步骤，创建子任务并分配给各模块，最后综合结果并返回给用户。其特点是把传统的计算逻辑和系统工具与大语言模型结合，通过提示引导模型使用这些工具。

例如，设想要对市售耳机进行市场调研，再根据价格、功能和用户评价选出几款。只凭预训练时接触的数据，原书讨论的基础 ChatGPT 无法完成这项任务。大语言模型智能体则可以读取请求，判断需要从网上查找信息，构造相关查询，执行搜索，下载结果网页，处理信息，最后给出若干建议。

不同类型的智能体一般都有若干共同模块，见图~\ref{fig:ch10-agent}。

\begin{figure}[!htbp]
\centering
\begin{tikzpicture}[
  font=\sffamily\small,
  module/.style={draw=blue!50!black,rounded corners=3pt,
    text width=2.35cm,minimum height=.67cm,align=center},
  tool/.style={module,text width=1.55cm,minimum height=.65cm,fill=green!14},
  exchange/.style={{Stealth[length=1.6mm]}-{Stealth[length=1.6mm]}},
  request/.style={-{Stealth[length=1.6mm]}}]
\node[module,fill=yellow!18] (core) at (0,0) {智能体核心};
\node[module,fill=blue!12] (plan) at (-3.85,0) {规划};
\node[module,fill=orange!16] (memory) at (3.85,0) {记忆};
\node[module,fill=black!11] (user) at (0,-1.72) {用户};
\node[tool] (web) at (-3.30,1.72) {网页搜索};
\node[tool] (calc) at (-1.10,1.72) {计算器};
\node[tool] (kb) at (1.10,1.72) {知识库};
\node[tool] (other) at (3.30,1.72) {其他工具};
\draw[exchange] (plan) -- (core);
\draw[exchange] (memory) -- (core);
\foreach \toolname in {web,calc,kb,other}{
  \draw[exchange] (core) -- (\toolname.south);
}
\draw[request] ([xshift=-.48cm]user.north) --
  node[left,font=\sffamily\scriptsize]{任务说明} ([xshift=-.48cm]core.south);
\draw[request] ([xshift=.48cm]core.south) --
  node[right,font=\sffamily\scriptsize]{汇总结果} ([xshift=.48cm]user.north);
\node[below,align=left,font=\sffamily\scriptsize] at (plan.south)
  {分解任务\\推理（思维链、Reflexion）};
\node[below,align=left,font=\sffamily\scriptsize] at (memory.south) {短期记忆\\长期记忆};
\draw (-4.17,2.25) -- (-.45,2.25) (0.45,2.25) -- (4.17,2.25);
\draw (-4.17,2.13) -- (-4.17,2.37) (4.17,2.13) -- (4.17,2.37);
\node at (0,2.25) {工具};
\end{tikzpicture}
\caption{大语言模型智能体的主要模块。核心接收用户任务，从记忆模块取得相关信息，交由规划模块分解子任务，再使用可用工具完成子任务。核心最后综合信息并将结果返回用户。}
\label{fig:ch10-agent}
\end{figure}


\begin{keypoint}
\begin{itemize}
\item \textbf{核心。}定义智能体的特性，接收用户指令，并向其他模块分配任务。这通过一个较长、详细的提示模板实现，用于指导大语言模型如何执行任务。
\item \textbf{规划。}确定完成整体任务所需的一系列步骤。采用思维链等推理方法（见第~4.5.4~节）以及 Reflexion 等迭代改进算法（\refcite{ch10-ref-022}{Shinn et al., 2023}），制定并完善一份由多个子任务组成的计划，再由智能体的各项功能执行。
\item \textbf{工具。}为智能体提供超出大语言模型本身能力的一系列功能。工具可以有很多种，例如网页搜索命令、代码编译器、计算器和各种 API 调用。
\item \textbf{知识。}智能体可在需要时查询的知识库。它可以类似第~7~章介绍的 RAG 系统，也可以是结构化数据库，由语言模型生成 SQL 等调用来查询。
\item \textbf{记忆。}记录从用户与智能体的交互中得到的信息，必要时供任务查阅。有时可分为短期记忆与长期记忆。短期记忆保留当前会话所有交互的细节，长期记忆则保存从多次交互中选取的相关信息。
\end{itemize}
\end{keypoint}

智能体可以针对科学写作（\refcite{ch10-ref-038}{Ziems et al., 2024}）、玩电子游戏（\refcite{ch10-ref-030}{Wang et al., 2023}）、操纵机器人（\refcite{ch10-ref-018}{Michael Ahn, 2022}）等具体任务设计。研究人员也开发了通用智能体，尝试执行用户提出的任意任务。早期例子是 AutoGPT。\footnote{\url{https://github.com/Significant-Gravitas/AutoGPT}}其结构与图~\ref{fig:ch10-agent}相近。它接收用户指令，用核心模块中设计好的提示模板组织工作流，通过思维链和自我批评拟定计划，再利用记忆模块与工具完成任务。它不接受用户对计划的反馈，而是自主尝试整个任务求解过程。另一种当时常用的智能体框架是 BabyAGI。\footnote{\url{https://github.com/yoheinakajima/babyagi}}其整体布局与 AutoGPT 相似，但每完成一个任务就会调整计划，而非依次执行预先确定的一串任务。

\begin{practicaltip}
许多智能体建立在为复杂大语言模型应用设计的开源软件包上。LangChain 和 LlamaIndex 是常见例子，见第~\ref{sec:ch08-frameworks}~节。这些软件包提供调用大语言模型的功能、常用工具的集成、适用于多种场景的提示模板以及 Web 托管功能。BabyAGI 的工作流使用了 LangChain 的集成组件。这种集成也是双向的，LangChain 将 AutoGPT 和 BabyAGI 纳入其产品，智能体系统可以方便地使用 LangChain 中已有的各种模型、向量索引和工具。
\end{practicaltip}

\subsection{大语言模型增强搜索}
\label{sec:ch10-search}
大语言模型的另一应用前沿，是增强不同环境下的搜索能力。大语言模型可用多种方式改进网页与文档搜索算法，大致分为以下几类。
\begin{itemize}
\item 改进排名靠前的搜索结果。
\item 改写或扩展查询。
\item 根据搜索结果进行推理。
\end{itemize}

\subsubsection*{改进排名靠前的搜索结果}
传统搜索根据查询做关键词匹配。它的一个主要限制是，无法返回主题相近却未包含相同关键词的文本。词向量扩大了单个词的搜索范围，使语义相近的词也能匹配。语言模型的句子嵌入进一步匹配语义相近的较长短语。随着嵌入模型的改善，检索有望更准确地限定相关范围，可用第~\ref{sec:ch07-reranking}~节介绍的重排序方法。先用高效搜索算法收集一组排名靠前的结果，再用较慢但能力更强的排序算法动态重排。以网页搜索为例，传统的关键词搜索可以返回几十万条匹配结果，初步排序较为粗略。可根据查询与网页内容的语义相似度，对前 $k$ 条结果重新排序。从 Transformer 出现以来，这类技术已在网页搜索行业普遍使用。能细致分析人类语言含义的大语言模型智能体，又扩展了这类方法的能力。

\subsubsection*{改进查询}
大语言模型还可用于改进查询。第~7~章在 RAG 背景下讨论过其中一些方法。
\begin{itemize}
\item \textbf{查询改写。}训练模型接收人的搜索查询，将其改写为更可能检索到相关结果的形式。
\item \textbf{查询到文档。}生成与所需搜索结果形式相似的合成文档，使之与目标文档更易匹配。
\item \textbf{查询到 SQL。}用大语言模型将针对结构化数据库的自然语言查询转换为代码调用，参见第~\ref{sec:ch07-structured}~节的 SQL 示例。\footnote{原书此处把该示例称为清单 7.4.1.1。第~7.4.1.1~节内包含相应示例。}
\end{itemize}

前瞻式主动检索增强生成（Forward-Looking Active REtrieval augmented generation，FLARE）也是提取更相关信息的一种技术（\refcite{ch10-ref-015}{Jiang et al., 2023b}）。在 FLARE 中，大语言模型从初始查询出发，设想可能包含相关信息的新上下文，接连生成后续查询，执行搜索并纳入新结果。

\subsubsection*{根据搜索结果进行推理}
大语言模型改变搜索的另一种方式，可类比第~7~章所述的检索增强生成。RAG 流水线本就包含对文档的语义搜索。基于 RAG 的搜索引擎可以只返回语义最相近的文档，供用户查阅。这种方法适用于在单份文档或网站域名内搜索，可以改善传统关键词搜索难以识别与查询相近但用词不完全相同的主题这一问题。

先搜索、再生成描述的 RAG 式方法也可用于网页搜索。大语言模型增强的搜索引擎可以在预先建立索引的网页中查找语义最接近的结果，提取页面信息并用作上下文，直接回答查询。You.com 创建了具有类似功能的网页入口。按原书当时的描述，You.com 采用与其他聊天机器人相似的网页查询界面，并与 Google 搜索结合。收到查询后，You.com 在网上搜索与问题有关的信息，处理网页的细节，再以这些上下文回答查询。由于能访问网上新近发布的信息，这类聊天机器人可将互联网作为资料库来做 RAG 式搜索。它可用于综合近期新闻中的信息、市场调研和一般问答。

这类应用仍有其他聊天机器人面临的问题，尤其是模型幻觉和信息来源的可靠性。一个值得注意的例子是，谷歌在原书写作时期开展人工智能新闻聚合与摘要试验，曾对简单问题给出十分不准确的回答，有时甚至让用户吃石头，或把胶水加到比萨上（\refcite{ch10-ref-010}{Grant, 2024}）。出现这些响应，是因为带有讽刺意味的评论或网页被用作上下文。由此可以看到，如果把互联网搜索结果等未经筛选的数据当成大语言模型问答的事实知识库，就可能发生错误。搜索并作答的系统一旦输出不准确的信息，可能造成伤害和声誉损失，因此开发时需要留意信息来源。

\section{结语}
\label{sec:ch10-closing}
在最后一章，我们试着把自己所看到的大语言模型创新前沿介绍给读者。更确切地说，这是多模态大语言模型的创新前沿。正如第~9~章所述，多模态正迅速成为语言建模的新范式。我们对前沿的看法来自撰写本书时查阅的大量文献，未来的面貌却难以预料。人类的创造与发明能力，始终是其中难以预测的变量。毕竟，当 Vaswani 等人发表奠基论文 \refcite{ch10-ref-029}{\textit{Attention Is All You Need}} 时，又有谁能预见 ChatGPT 后来展现的能力？

人类创新的不可预见性，使我们不敢对这一领域的未来作过于具体的判断。人类创新仍是一个影响因素，这本身也反映出大语言模型及其能力在何种程度上接近人类智能。我们要创造能与人类个体或群体同样胜任任务，甚至做得更好的机器，仍有许多东西要学习与理解，还会经历失败，也会取得更多进展。

我们希望，本书的内容和汇集的研究，能为有技能、兴趣、才智或机会参与这段人类技术发展进程的人提供参考。更重要的是，我们希望他们在推进技术时承担责任、遵循伦理，并充分尊重人的尊严。技术都有两面性，可能在众多任务中超越人类能力的技术尤其如此。

愿读者喜欢这段旅程，也希望本书能帮助你理解和掌握大语言模型的研究与应用。人工智能仍在发展，本书也需要随之更新。我们期待将来的版本，也希望读者同样期待。
